\chapter{Conclusion and Future Work}

%==============================================================================
\section{Conclusion}
%==============================================================================

This project set out to investigate whether tree-based flow matching, a generative framework originally designed for static tabular data, could be extended to synthesise realistic sequential ICU patient trajectories from the MIMIC-III clinical database. The resulting system implements a two-stage autoregressive pipeline with two interchangeable model backbones (ForestFlow and HS3F), a multi-axis evaluation stack, and a reproducible hyperparameter sweep infrastructure. The key findings and contributions are summarised below, organised by the research questions posed in Chapter~1.

\subsection{Research Question 1: Feasibility of Autoregressive Extension}

The project demonstrates that tree-based flow matching can be adapted to sequential clinical data through autoregressive conditioning. By factorising the joint trajectory distribution into an IID hour-0 model and a conditional transition model, where each hourly state is generated given lagged dynamics and static covariates, the system produces multi-hour patient trajectories without requiring recurrent or attention-based neural architectures. The implementation trains entirely on CPUs using XGBoost regressors at discrete time levels, confirming the computational efficiency advantage of tree-based approaches.

The hyperparameter sweep covers a Fibonacci-scale $8 \times 8 = 64$-cell grid over $n_t, n_{\text{noise}} \in \{1,2,3,5,8,13,21,34\}$; $62$ of $64$ cells are complete in the canonical CSV snapshot used for this draft, with the remaining unfilled cells at $(n_t, n_{\text{noise}})=(34,21)$ and $(34,34)$. The sweep is interpreted primarily as an effect-analysis study of how utility, quality, privacy, and compute respond to $n_t$ and $n_{\text{noise}}$, with ``best observed'' settings treated only as summaries of the tested region.
Across completed cells, clinical range violation rates drop from $\approx 72\%$ at low $n_t$ to a best of $\approx 52\%$ at $(n_t{=}21, n_{\text{noise}}{=}2)$, and $n_t$ is the dominant hyperparameter on every utility and quality axis while $n_{\text{noise}}$ has a near-flat effect.

Feasibility is established: mortality ROC-AUC increases with $n_t$ and then visibly saturates, while $n_{\text{noise}}$ is largely flat on utility and quality but increases compute. The best observed mean cell reaches ROC-AUC $0.658$ at $(n_t{=}21, n_{\text{noise}}{=}1)$ against a real-model baseline of $\approx 0.82$, but this value is interpreted as a descriptive in-grid maximum rather than an optimisation endpoint. A saturating-fit extrapolation to the production configuration $(n_t{=}50, n_{\text{noise}}{=}100)$ suggests a mortality ROC-AUC plateau near $0.65 \pm 0.04$.

This extrapolation is anchored on three-seed-per-cell means from a sparse Fibonacci grid with two unfilled cells in the high-$n_t$ corner and no formal model-based confidence interval on the best cell; the three candidate saturating forms agree on the populated grid but diverge outside it, so the value should be read as indicative rather than as an empirical estimate. Under all three forms the extrapolation does not close the gap to the real-model baseline, which is consistent with the view that the residual gap is structural rather than hyperparameter-limited and that further architectural improvements are needed for clinical-grade synthesis; empirical validation at $(50, 100)$ remains necessary before this interpretation can be treated as established. The dissertation claim is therefore effect-level (trend shape, saturation, and trade-offs), not prediction-level certainty outside the observed grid.

\subsection{Research Question 2: Heterogeneous Feature Handling}

Two model backbones were implemented to address this question. ForestFlow treats all target dimensions, including binary intervention flags, as continuous flows, while HS3F generates features sequentially, routing continuous variables through flow-matching regressors and categorical variables through XGBoost classifiers. The HS3F approach is architecturally better aligned with the mixed-type nature of ICU data, where binary treatment indicators (ventilation, vasopressors) and categorical variables coexist with continuous vital signs and laboratory measurements.

Backbone-axis evidence at this stage is a three-seed matched-cell comparison at $(n_t{=}5, n_{\text{noise}}{=}5)$ that directly contrasts HS3F, ForestFlow, and the CTGAN baseline on identical training rows, synthetic budgets, and trajectory seeds (see Evaluation, section~\ref{sec:backbone-comparison}). The HS3F full-grid sweep is in progress; the ForestFlow and CTGAN full grids are not run and are tracked as an active next step in the project's status log.

The matched cell shows that under balanced accuracy, the appropriate mortality metric given the ICU class imbalance, all three backbones cluster near chance (HS3F $0.501$, CTGAN $0.507$, ForestFlow $0.535$) with overlapping three-seed confidence intervals, so no backbone produces a downstream classifier that reliably discriminates mortality at a fixed threshold. On ROC-AUC, which is invariant to class prevalence and measures ranking quality, HS3F leads at $0.672$ against ForestFlow's $0.619$ and CTGAN's $0.497$: the HS3F--CTGAN gap is statistically significant, whereas the HS3F--ForestFlow gap is not significant under the three-seed budget (CI95 half-widths $\approx 0.021$ and $\approx 0.077$ respectively).

On trajectory-level temporal coherence the matched cell separates the two tree-based backbones more cleanly than patient-level TSTR does: HS3F's autocorrelation distance ($0.22$) is roughly half of ForestFlow's ($0.44$), and its temporal cross-feature correlation drift ($0.26$) is likewise well below ForestFlow's ($0.41$), consistent with the HS3F design hypothesis that explicit categorical heads stabilise the autoregressive conditioning chain. HS3F also trains $\approx 32\%$ faster than ForestFlow at this cell ($2194$~s vs.\ $3224$~s) and $\approx 3.7\times$ faster than CTGAN. Taken together, ROC-AUC leadership, roughly 2$\times$ advantage on autocorrelation distance and $\approx 1.6\times$ on cross-feature drift, privacy parity, and the lowest wall-clock cost, these signals motivate HS3F as the primary backbone used for the main sweep grid, with the caveat that the single matched cell and three-seed budget make the HS3F-over-ForestFlow ordering indicative rather than definitive.

\subsection{Research Question 3: Utility and Privacy}

The expanded TSTR evaluation reveals task-dependent utility transfer with a clear diagnostic pattern. Mortality F1 remains in $0.04$--$0.23$ across the completed grid versus a real-model F1 of $\approx 0.47$, and LOS macro-F1 stays in $0.17$--$0.38$ versus a real-model macro-F1 of $\approx 0.77$. This gap confirms that the generator preserves majority-class ranking but collapses minority-class support, a class-collapse failure mode rather than a total distributional failure (detailed in Evaluation section~\ref{sec:mortality-tstr}). The best observed mortality ROC-AUC mean ($0.658$ at $(n_t{=}21, n_{\text{noise}}{=}1)$) closes about $49\%$ of the gap from chance ($0.5$) to the real-model baseline ($\approx 0.82$), indicating meaningful but incomplete utility transfer.

Privacy evaluation infrastructure is implemented and integrated into the sweep pipeline: DCR summary statistics (median, low-tail quantiles, exact-match rate), baseline and overfitting protection scores following SDMetrics formulations, and a DOMIAS-inspired membership inference attack reporting ROC-AUC and attacker advantage. The combination of proximity-based and attack-based metrics follows current best practice for synthetic health data \cite{kaabachi2025scoping,breugel2023domias}. The implementation explicitly cautions that DCR-style proximity metrics are useful diagnostics but do not constitute standalone privacy guarantees, a finding emphasised in recent work showing distance-only proxies can miss true membership leakage \cite{yao2025dcrdelusion}. Across all completed sweep cells, the DCR baseline-protection score saturates at $1.0$ under the SDMetrics $\min(1,\cdot)$ formulation, no synthetic record exactly matches a training record (\code{dcr\_exact\_match\_rate} $=0$ throughout), and the DOMIAS-style MIA ROC-AUC is confined to the narrow range $0.594$--$0.599$, above chance but well below the bands typically cited for clear leakage. Under the measured configurations the generator therefore does not exhibit gross memorisation. However, even with three-seed aggregation across the main HS3F grid, the attacker remains a nearest-neighbour DOMIAS approximation rather than a full shadow-model attack, and the protection score saturates at the SDMetrics ceiling, so these results should be interpreted as diagnostic indicators rather than definitive privacy bounds. Establishing formal release thresholds requires a larger-scale attack battery, which is deferred to future work.

\subsection{Final Judgement}

The final judgement is:
\begin{itemize}\tightlist
    \item \textbf{Successfully achieved:} a complete research system implementation that extends tree-based flow matching to sequential ICU trajectory synthesis and evaluates it across utility, fidelity, privacy diagnostics, and CPU efficiency.
    \item \textbf{Partially achieved:} measurable downstream utility transfer, mainly on ranking metrics, with clear but incomplete gains over chance and over the CTGAN lower-bound baseline.
    \item \textbf{Not achieved:} practical release readiness. The observed utility gap to real-data baselines, persistent minority-class collapse, and high range-violation rates show the system is not yet adequate for deployment-oriented data release.
    \item \textbf{Not established:} formal privacy release assurance. Current privacy evidence is diagnostic and encouraging in parts, but insufficient for a formal release claim.
\end{itemize}

\subsection{Summary of Contributions}

This project is an application-level contribution: it does not propose new generative algorithms but integrates existing methods into a new domain and problem structure. The principal contributions are:
\begin{enumerate}\tightlist
    \item \textbf{Domain adaptation:} A two-stage autoregressive pipeline extending tree-based flow matching from static tabular benchmarks to sequential ICU trajectory synthesis on MIMIC-III, a setting not previously explored with this method family.
    \item \textbf{Dual-backbone implementation:} Both ForestFlow and HS3F backbones implemented from the original papers with a unified training, generation, and evaluation interface, enabling backbone comparison on identical clinical data.
    \item \textbf{Multi-axis evaluation framework:} A single integrated evaluation pipeline that combines multi-task patient-level TSTR (mortality, ICU length-of-stay), distributional quality metrics (KS, correlation, clinical ranges), trajectory-level temporal coherence metrics (autocorrelation, transition smoothness, cross-feature drift), and privacy-risk diagnostics (DCR, DOMIAS-style membership inference), so that every sweep cell produces utility, fidelity, temporal, and privacy measurements on the same synthetic cohort.
    \item \textbf{Reproducible experiment infrastructure:} Hyperparameter sweep with caching, resumability, degeneracy detection, and canonical result logging, supporting systematic investigation of configuration effects.
    \item \textbf{Honest empirical analysis:} Application to MIMIC-III ICU data with transparent reporting of both strengths (correlation preservation, range violation reduction at high $n_t$) and limitations (F1 collapse, structural utility gap, single-dataset scope).
\end{enumerate}

%==============================================================================
\section{Limitations}
%==============================================================================

The following limitations bound the conclusions of this work:
\begin{itemize}\tightlist
    \item \textbf{Single dataset:} All experiments use MIMIC-III via a MIMIC-Extract style pipeline. Generalisability to other hospitals, charting systems, or patient populations is untested.
    \item \textbf{Asymmetric baseline capability:} The CTGAN baseline is a lower-bound reference rather than a fully matched autoregressive competitor (see Evaluation, section~\ref{sec:backbone-comparison}).
    \item \textbf{Limited seed count per sweep cell:} The main HS3F grid is aggregated over three trajectory seeds per cell. This improves robustness over single-seed reporting, but remains a low-shot uncertainty budget for tight statistical separation between close configurations.
    \item \textbf{Persistent range violations:} Although range violation rates drop to a best of $\approx 52\%$ at $(n_t{=}21, n_{\text{noise}}{=}2)$ (from $\approx 72\%$ at low $n_t$), this remains far from a deployable threshold for clinical applications and indicates that marginal distributions are not yet fully aligned.
    \item \textbf{Fixed-length trajectory generation:} The current 48-hour rollout design does not match real ICU stay-length variability, which limits LOS utility and leaves stay-length KS structurally poor (evidence in Evaluation section~\ref{sec:sweep_quality}).
    \item \textbf{Speculative extrapolation to the production configuration:} Predictions at $(n_t{=}50, n_{\text{noise}}{=}100)$ are model-based trend estimates on a partially populated grid and should not be read as empirical measurements.
    \item \textbf{Limited temporal coherence evaluation:} Basic trajectory-level metrics (autocorrelation distance, transition smoothness, stay-length distribution, cross-feature temporal drift) have been implemented, but more sophisticated temporal assessments, such as dynamic time warping, Granger causality, or learned discriminators on trajectory windows, remain unexplored.
\end{itemize}

%==============================================================================
\section{Future Work}
%==============================================================================

Several directions would strengthen and extend this work:

\paragraph{ForestFlow backbone comparison.}
The current sweep covers HS3F only. Running the same grid with the ForestFlow backbone would directly answer Research Question~2 and allow empirical comparison of all-continuous versus heterogeneous sequential generation on ICU data. Increasing the number of trajectory seeds and adding bootstrap confidence intervals would further strengthen conclusion validity.

\paragraph{Baseline architecture parity.}
Future work should add a truly autoregressive baseline with explicit per-row conditional generation (for example, a sequential GAN variant or a redesigned conditional wrapper around CTGAN-like components); the architectural and scope rationale is set out in Evaluation section~\ref{sec:baseline-scope}. Implementing this well requires expanded validation and additional experimental budget beyond this thesis, which would enable stronger parity claims in backbone comparisons.

\paragraph{Additional TSTR tasks.}
Extending the TSTR framework with further clinical prediction tasks, such as readmission prediction or specific intervention timing, would provide a more robust assessment of downstream utility across diverse clinical endpoints.

\paragraph{Advanced temporal coherence metrics.}
The current trajectory-level metrics (autocorrelation distance, transition smoothness, stay-length KS, cross-feature temporal drift) provide a first-order assessment of temporal plausibility. More sophisticated approaches, such as dynamic time warping for trajectory similarity, Granger causality tests for lead--lag relationships between clinical variables, or learned trajectory discriminators, would provide deeper insight into temporal fidelity and could guide architectural improvements to the autoregressive generation scheme.

\paragraph{Constraint enforcement and post-processing.}
Incorporating clinical range clipping, logical constraint checking (e.g., sex--procedure consistency), or learned post-processing steps could reduce the high range violation rates observed in the current results.

\paragraph{Subgroup fairness evaluation.}
Stratifying evaluation metrics by demographic subgroups (age, sex, ethnicity) would reveal whether the generator disproportionately degrades fidelity or utility for minority populations, addressing a key ethical concern in synthetic health data.

\paragraph{Cross-cohort replication.}
Applying the pipeline to a second clinical dataset (e.g., eICU-CRD or a UK hospital system) would test external validity and the portability of the two-stage architecture.

\paragraph{Privacy hardening.}
Running a full battery of membership inference attacks with varied attacker knowledge assumptions, and establishing release thresholds for synthetic data sharing, would strengthen the privacy claims beyond the current diagnostic-level assessment.
